\begin{table}[!ht]
\centering
\caption{\label{tab:country-list}Country list.}
\centering
\begin{threeparttable}
\begin{tabular}[t]{llllll}
\hline
\multicolumn{3}{c}{High-income EU} & \multicolumn{3}{c}{Middle-income EU} \\
\cline{1-3} \cline{4-6}
Country & N of observations & N of affiliates & Country & N of observations & N of affiliates\\
\hline
Austria & 144 & 18 & Bulgaria & 24 & 3\\
Belgium & 760 & 95 & Czechia & 400 & 50\\
Denmark & 128 & 16 & Greece & 48 & 6\\
Finland & 88 & 11 & Hungary & 256 & 32\\
France & 1584 & 198 & Poland & 280 & 35\\
Germany & 3520 & 440 & Portugal & 112 & 14\\
Ireland & 168 & 21 & Romania & 72 & 9\\
Italy & 752 & 94 & Slovakia & 80 & 10\\
Luxembourg & 80 & 10 & Slovenia & 24 & 3\\
Netherlands & 2128 & 266 & Spain & 576 & 72\\
Sweden & 232 & 29 &  &  & \\
United Kingdom & 3968 & 496 &  &  & \\
\hline
\end{tabular}
\begin{tablenotes}
\item \textit{Note: } 
\item This table shows the number of observations ($=$ the number of affiliates in 2010 $\times$ eight years) by the host country. The sample includes the Japanese affiliates in the EU. Sample restricted to affiliates with $>10$\% Japanese ownership.
\end{tablenotes}
\end{threeparttable}
\end{table}
